\chapter{La idea en pocas páginas}
\label{cap:s-idea}

\section{El problema: un estado lleno de un lado y vacío del otro}
Quintana Roo tiene el Caribe más visitado de México, pero los visitantes no se reparten parejo. Casi todos llegan al
norte: Cancún, la Riviera Maya y Tulum. El sur, con Chetumal, la bahía, las pirámides en la selva y los pueblos mayas,
recibe muy poca gente aunque tiene espacio.

En 2026 el desequilibrio se volvió un problema visible:
\begin{itemize}
\item \textbf{Sargazo récord}: más de 104,700 toneladas recogidas y 56 de 140 playas del estado en rojo en agosto.
\item \textbf{Tulum en crisis}: las ventas cayeron hasta 60\,\% y cerraron negocios. Las visitas a su zona arqueológica
  bajaron 31.3\,\% de enero a julio de 2026 contra el mismo periodo de 2025 (INAH).
\item \textbf{Cozumel saturado por cruceros}, Isla Mujeres con carencias de agua y drenaje, Holbox con basura acumulada.
\end{itemize}

El \emph{Problema Prototípico} de la UNRC pide justo esto: diseñar una campaña de publicidad basada en datos que atraiga
visitantes a un destino \textbf{de forma responsable}, que reparta mejor a los turistas y que aproveche el espacio que
sobra sin crear una nueva saturación.

\section{La cifra que lo resume todo}
\begin{ejemplo}[ · cuánto tiene el sur y cuánto recibe]
Los cinco lugares de la campaña tienen:
\begin{itemize}
\item el \textbf{12.3\,\%} de la gente del estado (Censo 2020) y el \textbf{14.4\,\%} de sus negocios turísticos (DENUE);
\item pero solo el \textbf{1.4\,\%} de los pasajeros de avión, el \textbf{1.7\,\%} de los cuartos de hotel y el
  \textbf{4.1\,\%} de los visitantes a zonas arqueológicas.
\end{itemize}
En palabras sencillas: \textbf{los negocios siguen a la gente que vive ahí, pero los turistas no}. El sur tiene la
infraestructura de una región que podría recibir más visitantes, y no los recibe.
\end{ejemplo}

Otra forma de verlo: en 2025, la zona arqueológica de Tulum recibió \textbf{1,031,443} visitantes; las tres zonas de la
Ruta de las pirámides del sur (Kohunlich, Dzibanché e Ichkabal) recibieron \textbf{66,628}. Tulum recibe
\textbf{15.5 veces} más gente. Y en Chetumal, en 2024, de cada 10 cuartos de hotel solo se ocuparon unos 6: se usaron
458,696 noches de 791,016 disponibles, el 58.0\,\%.

\section{Los cinco lugares de la campaña}
\begin{center}
\small
\begin{tabularx}{\linewidth}{>{\raggedright\arraybackslash}p{0.25\linewidth}Y}
\encabezado \h{Lugar} & \h{Qué es y su cifra principal} \\
\textbf{Chetumal} & La capital del estado, frente a la bahía, con el Museo de la Cultura Maya. 169,028 habitantes. En
2024 tuvo 58.0\,\% de ocupación hotelera: 4 de cada 10 cuartos vacíos. Recibe gente por la frontera con Belice (49,097
cruces en junio de 2026) y por el Tren Maya. \\
\filagris\textbf{Bahía de Calderitas y Oxtankah} & Un pueblo de pescadores a 7.7 km de Chetumal y la zona arqueológica de
Oxtankah, la ciudad maya más grande de la bahía. Oxtankah recibió 11,017 visitantes en 2025, el 94.8\,\% mexicanos. \\
\textbf{Ruta de las pirámides} & Kohunlich, Dzibanché e Ichkabal, tierra adentro y sin sargazo. 66,628 visitantes en
2025, unas 183 personas al día entre las tres. \\
\filagris\textbf{Maya Ka'an y Kantemó} & Felipe Carrillo Puerto, Tihosuco, Señor, Chunhuhub y Kantemó: turismo
comunitario maya, con estación del Tren Maya. 44,388 habitantes. \\
\textbf{Laguna Milagros y Xul-Ha} & Una laguna dulce cerca de Chetumal, del mismo sistema que Bacalar. Es el lugar con
menos datos (ninguna serie turística propia), así que se promueve con límite estricto. \\
\bottomrule
\end{tabularx}
\normalsize
\end{center}

\textbf{¿Y el norte?} Cancún, la Riviera Maya y Tulum \textbf{no se promueven}. Aparecen solo como \emph{referencia}: son
el punto de comparación (sin ellos no se puede medir que el sur ``tiene espacio'') y son de donde viene buena parte del
público de la campaña. En la página, Cancún y la Riviera Maya salen en el planeador con la etiqueta ``Referencia: la
campaña no lo promueve'', y si su mes está lleno, la página recomienda un lugar del sur.

\begin{porque}
Al principio se eligieron 8 lugares. Brandon pidió quitar todo lo que tuviera sargazo, cierres o relación con Tulum, y
salieron tres: \textbf{Cobá} (está en el municipio de Tulum), \textbf{Muyil} (cerrada de junio de 2024 a febrero de 2026 y
pegada a Tulum) y \textbf{la Ribera del Río Hondo} (ninguna fuente oficial publica sus visitantes). También se descartaron
Chacchoben (el 92\,\% de sus visitantes llega en crucero por Mahahual, que tiene sargazo) y Bacalar (su laguna está en
deterioro ecológico).
\end{porque}

\section{La respuesta: una campaña con tres pantallas}
La propuesta se llama \textbf{``El sur tiene espacio''}. Invita a la gente al sur, pero con reglas de cuidado: nunca
anuncia un lugar en su temporada alta, se pausa sola si hay mal clima o si llega más gente de la esperada, y a quien iba
al norte lleno le ofrece el sur.

Para tomar esas decisiones, el sistema funciona como la torre de control de un aeropuerto, con tres pantallas que
responden tres preguntas distintas y nunca se enciman:

\begin{center}
\small
\begin{tabularx}{\linewidth}{lYY}
\encabezado \h{Pantalla} & \h{Pregunta} & \h{Ejemplo de respuesta real} \\
\textbf{El Radar} & ¿Dónde hay espacio \emph{hoy}? & En julio de 2026 los cinco lugares del sur están tranquilos;
Cancún, concurrido. \\
\filagris\textbf{El Pronóstico} & ¿\emph{Cuándo} conviene ir en los próximos 12 meses? & En enero la Ruta recibe 61\,\% más
gente que un mes normal; en septiembre, la mitad. Mejor ir en noviembre. \\
\textbf{La Torre} & ¿Qué hace la campaña \emph{esta semana}? & Esta semana pasa una tormenta: el anuncio de ese lugar se
pausa y el dinero se guarda para la siguiente semana buena. \\
\bottomrule
\end{tabularx}
\normalsize
\end{center}

\begin{intuicion}
Piensa en un viaje en carretera. El Radar es mirar por la ventana: cómo está el tráfico ahora. El Pronóstico es revisar
el calendario de puentes y vacaciones: cuándo va a haber tráfico. La Torre es el copiloto que, en el camino, dice ``hay
un choque adelante, mejor tomemos otra ruta''.
\end{intuicion}

\section{Cómo se pasan la información}
Las tres pantallas trabajan en cadena:
\begin{enumerate}
\item \textbf{El Pronóstico arma el plan}: cuántos visitantes se esperan cada mes y en qué meses hay riesgo de llenarse.
  Con eso se decide cuánto dinero va a cada lugar, cada mes y cada red social.
\item \textbf{El Radar dice el estado}: tranquilo, concurrido o saturado, y qué tan probable es que cambie.
\item \textbf{La Torre ejecuta}: cada semana compara lo que pasó con lo esperado y, si hace falta, pausa un anuncio o
  enciende otro.
\item \textbf{La campaña} pone el mensaje, el público y el canal de cada anuncio.
\end{enumerate}

\section{Cómo se trabajó: 12 fases, una a la vez}
El proyecto se hizo en 12 fases (0 a 11). No se abría una fase sin terminar y revisar la anterior.

\begin{center}
\small
\begin{tabularx}{\linewidth}{rlY}
\encabezado \h{Fase} & \h{Nombre} & \h{Qué se hizo} \\
0 & Cimientos & El plan, las reglas y la computadora lista para procesar millones de datos \\
\filagris 1 & Descarga & 353 archivos oficiales, 8,134,802 registros \\
2 & Limpieza & Los datos limpios, ordenados y consultables (56 tablas documentadas) \\
\filagris 3 & Planteamiento & Elegir y medir los 5 lugares; medir la desigualdad \\
4 & Radar & Índice de presión, estados y predicción del mes siguiente \\
\filagris 5 & Pronóstico & Visitantes de los próximos 12 meses, tormentas y escenarios \\
6 & Presupuesto & El reparto óptimo del dinero de la campaña \\
\filagris 7 & Torre en vivo & La campaña semana a semana \\
8 & Campaña & Viajeras ideales, marca, mensajes, medios e indicadores \\
\filagris 9 & Servidor & La página calcula en vivo con los modelos \\
10 & Página final & Revisión de accesibilidad \\
\filagris 11 & Cierre & Trazabilidad, mapa del informe y guion del coloquio \\
\bottomrule
\end{tabularx}
\normalsize
\end{center}
